\section{Introduction}
\label{sec:introduction}

The Android operating system accounts for the overwhelming majority of
mobile devices in use worldwide, and consequently for the overwhelming
majority of mobile malware samples captured in
the wild~\cite{alswaina2020android,kaspersky2023mobile}. Attackers
exploit both the openness of the ecosystem and the relative permissiveness
of the permission model to deploy families that engage in SMS fraud,
banking credential theft, device bricking, click fraud, and ransomware.
Because modern malware families are repackaged and minor-version-bumped
hundreds of times, \emph{family-level classification} — assigning a
new sample to a known family — has become more operationally useful
than hash-level detection.

\subsection{The static--dynamic trade-off}

Static analysis is fast but increasingly brittle: packers, control-flow
flattening, and Java reflection collapse the observable API surface to
generic entry points, so static features can no longer distinguish
family-defining behaviour. Dynamic analysis — executing samples inside an
instrumented Android sandbox — is slower but observes the \emph{true}
behaviour. Hooking frameworks such as Droidmon emit a
time-ordered log of every framework-level API call. These logs form the
input to our study.

\subsection{Limitations of existing approaches}

Two lines of prior work dominate.

\paragraph{Statistical rule-mining} D'Angelo et al.~\cite{dangelo2023association}
treat each sample as a $k$-spaced Markov chain of API transitions
$\{A \rightarrow B\}$, mine associative rules with
support/confidence pruning, and classify new samples via
rule-voting. This approach is interpretable and fast, but it (i) loses
long-range sequential context by only modelling local transitions,
(ii) has no learned representation, and (iii) is sensitive to pruning
thresholds.

\paragraph{Deep attention models} Transformer-based classifiers
produce strong results on API sequences, but standard softmax attention
produces ``attention sinks''~\cite{qiu2025gated} and offers no intrinsic
mechanism for per-sample explanation. Post-hoc methods such as
SHAP or LIME can be applied, but their explanations are approximations
of a surrogate model, not of the classifier itself.

\subsection{This paper}

We propose GAME-Mal (\textbf{G}ated \textbf{A}ttention over
\textbf{M}arkov \textbf{E}mbeddings for \textbf{Mal}ware
classification). Our contributions are:

\begin{itemize}
  \item A \textbf{reflection-aware API resolution} scheme that uses
  \texttt{hooked\_class.hooked\_method} for reflective calls (preserving
  the true target behind \texttt{Method.invoke}) while retaining a
  \texttt{REFL:} prefix as a mechanism tag.

  \item The first application, to our knowledge, of
  \textbf{sigmoid-gated multi-head attention} (Qiu et al.\
  G1 position) to Android malware classification. A head-specific
  sigmoid gate after the scaled dot-product attention output introduces
  non-linearity, input-dependent sparsity, and breaks the linear
  bottleneck between $V$ and $W_O$.

  \item An \textbf{intrinsic explainability} mechanism: the gate
  activations themselves rank APIs by importance per sample, yielding
  per-family behavioural fingerprints without any post-hoc method.

  \item A \textbf{reproducible experimental pipeline} that runs
  four sklearn baselines (Random Forest, Linear SVM, Decision Tree,
  Gaussian Naïve Bayes), a faithful reproduction of
  D'Angelo et al.'s rule-voting classifier, and GAME-Mal under
  identical stratified $k$-fold splits.
\end{itemize}

The remainder of the paper is organised as follows.
Section~\ref{sec:related} reviews prior work on Android malware
classification, Markov chain rule mining, and gated attention.
Section~\ref{sec:methodology} details the GAME-Mal architecture and
the reflection-aware preprocessing pipeline.
Section~\ref{sec:experiments} describes the dataset and experimental
protocol. Section~\ref{sec:results} reports quantitative results.
Section~\ref{sec:explainability} analyses the gate activations and
extracts per-family API fingerprints.
Section~\ref{sec:discussion} discusses limitations.
Section~\ref{sec:conclusion} concludes.
